






\begingroup
\setslidetemplate
\begin{frame}[t]
  \frametitle{Additive Log-Ratio (ALR) Transformation}

  \vspace{0.2cm}
\note{One of the earliest and most influential ideas was to transform the data from the simplex into real space using log-ratio transformations. Introduced By Aitchison (1982). 

One of the simplest examples is the additive log-ratio transformation, or ALR.

In this approach, we choose one component as a reference, divide the remaining components by that reference, and then take logarithms.

This maps the data into real space, where standard statistical methods become easier to apply, and the covariance matrix is full rank.

However, this comes with a limitation: the result depends on the choice of the reference component.

If we change the reference, the transformed data also change.}

  \sechead{Definition.}
  For $\mathbf{x}=(x_1,\ldots,x_K)^\top \in \mathcal{S}^{K-1}$,
  choose a reference component $x_K$:

  \vspace{0.35cm}
  \[
    \mathrm{ALR}(\mathbf{x})
    =\left(
      \log\frac{x_1}{x_K},\;
      \log\frac{x_2}{x_K},\;\ldots,\;
      \log\frac{x_{K-1}}{x_K}
    \right)
    \in \mathbb{R}^{K-1}
  \]

  \vspace{0.3cm}

  \only<1>{
    \begin{itemize}\setlength\itemsep{6pt}
      \item Maps compositional data to Euclidean space $\mathbb{R}^{K-1}$
      \item Covariance matrix becomes full rank $\Rightarrow$ standard methods applicable
      \item \textcolor{red!70!black}{Depends on choice of reference component $x_K$}
    \end{itemize}

    \vspace{0.3cm}
    \refcite{Aitchison (1982)}
  }

\end{frame}
\endgroup
 
% ══════════════════════════════════════════════════════════════
% SLIDE 2 — CLR
% ══════════════════════════════════════════════════════════════
\begingroup
\setslidetemplate
\begin{frame}[t]
  \frametitle{Centered Log-Ratio (CLR) Transformation}

  \vspace{0.2cm}
\note{To avoid this dependence, Aitchison (1982) introduced another transformation the centered log-ratio transformation, or CLR.

Here we divide each component by the geometric mean of all components.


But,  the covariance matrix remains singular, which leads us to the next transformation.}
  \sechead{Definition.}
  Normalize by the geometric mean:

  \vspace{0.3cm}
  \[
    \mathrm{CLR}(\mathbf{x})
    = \left(
      \log\frac{x_1}{g(\mathbf{x})},\ldots,
      \log\frac{x_K}{g(\mathbf{x})}
    \right)
    \in \mathbb{R}^K
  \]

  \[
    g(\mathbf{x}) = \left(\prod_{j=1}^K x_j\right)^{1/K}
  \]

  \vspace{0.2cm}

  \only<1>{
    \begin{itemize}\setlength\itemsep{6pt}
      \item Symmetric representation (no reference component)
      \item Components satisfy: $\sum_{i=1}^K \mathrm{CLR}_i = 0$
      \item Still lies in a constrained subspace $\Rightarrow$ singular covariance
    \end{itemize}

    \vspace{0.3cm}
    \refcite{Aitchison (1982)}
  }

\end{frame}
\endgroup
 
% ══════════════════════════════════════════════════════════════
% SLIDE 3 — ILR
% ══════════════════════════════════════════════════════════════
\begingroup
\setslidetemplate
\begin{frame}[t]
  \frametitle{Isometric Log-Ratio (ILR) Transformation}
  \vspace{0.2cm}
\note{To address this, one more step is added.

By projecting the CLR coordinates onto an orthonormal basis, usually with a Helmert sub-matrix, we obtain the isometric log-ratio transformation, or ILR.

This gives a representation in real Euclidean space with full rank covariance.

Among these methods, ILR is one of the most commonly used in practice.}
 
      \sechead{Definition.} (\refcite{Barceló-Vidal et al.\ (2001); Egozcue et al.\ (2003)})

      By projecting the CLR coordinates onto an orthonormal basis, usually with a $(K{-}1)\!\times\! K$ Helmert sub-matrix $\mathbf{H}$, we obtain the isometric log-ratio transformation, or ILR:
 
      \vspace{0.2cm}
      \[
        \mathrm{ILR}(\mathbf{x})
        = \mathbf{H}\;\mathrm{CLR}(\mathbf{x})
        = (z_1,\;z_2,\;\ldots,\;z_{K-1})
        \in\mathbb{R}^{K-1}
      \]
      \[
        \mathbf{H}\mathbf{H}^\top = \mathbf{I}_{K-1}
        \quad\text{(rows orthonormal)}
      \]
 \only<1>{
      \begin{alertblock}{Properties.}
      \begin{itemize}\setlength\itemsep{4pt}
        \item Resolves CLR singularity: $\Sigma_{\mathrm{ILR}}$ is \textbf{full rank}
        \item Any $(K{-}1)\!\times\!K$ orthonormal matrix works in place of $\mathbf{H}$
        \item \textcolor[HTML]{2A9D8F}{\textbf{Isometric:}} Distances preserved by $\mathbf{H}$:
          $\|\mathrm{CLR}(\mathbf{x})-\mathrm{CLR}(\mathbf{y})\| = \|\mathrm{ILR}(\mathbf{x})-\mathrm{ILR}(\mathbf{y})\|$
       
        
      \end{itemize}
      
            \end{alertblock}

 }

\end{frame}
\endgroup



\begin{frame}{One Idea: Transforming to Real Space}


\only<1>{
\note{These transformations are useful, and in many applications they work very well.

But A group of researchers strongly argued that log-ratio transformations are not just useful, but necessary for compositional data analysis.

To support this view, they introduced a set of principles that any valid method for compositional data should satisfy.


}
\begin{itemize}
    \item[1)]
    CoDA School: Some researchers claim that log-ratio transformations are necessary for compositional data.
    \item[2)]
    To support this view, a set of principles (CoDA principles) was proposed that any valid compositional method should satisfy.
\end{itemize}

  \vspace{0.2cm}
  \renewcommand{\arraystretch}{1.35}
  \setlength{\tabcolsep}{8pt}

  \begin{tabular}{>{\bfseries}l l >{\footnotesize\itshape\color{gray}}l}
    \toprule
    \rowcolor{mainblue}
    \color{white}Method &
    \color{white}Formula &
    \color{white}Reference \\
    \midrule

    \rowcolor{accentblue!20}
    ALR & $\log x_i - \log x_K$ & Aitchison (1982) \\

    \rowcolor{accentblue!20}
    CLR & $\log x_i - K^{-1}\sum_{j=1}^{K} \log y_j$ & Aitchison (1983) \\

    \rowcolor{accentblue!20}
    ILR & $\mathbf{H}\!\left(\log y_i - K^{-1}\sum_{j=1}^{K} \log y_j\right)$ & Egozcue et al.\ (2003) \\

    \midrule

    \bottomrule
  \end{tabular}

}



\only<2>{

     \note{
The three most well-known are scale invariance, permutation invariance, and subcompositional coherence.

Scale invariance means that changing the scale of the composition should not affect the result.

Permutation invariance means that the analysis should not depend on the order of the components.

And subcompositional coherence means that conclusions based on a subset of components should remain consistent with conclusions from the full composition.

Based on these principles, researchers developed new definitions of distance, geometry, and statistical operations specifically for compositional data.


I mention these principles because they appear everywhere in the compositional data literature.



I spent more than enough time trying to understand why they are important?

And I failed. I don't know what role they play in statistical analysis of compositional data.

But I realized that I'm not alone in this and that there has been a lot of disagreement about these principles and logarithmic ratio transformations.

I would like to briefly mention one smal part of that debate before moving on.



}

   \begin{alertblock}{Principles of CODA}
      \begin{itemize}
       \item 
   \textbf{scale invariance},
   \item
   \textbf{permutation invariance}, and 
   \item \textbf{subcompositional coherence}
   If \( \Delta_p(\mathbf{x}, \mathbf{y}) \) denotes a distance measure between compositions with \( p \) parts, then
    \begin{equation*}
    \Delta_k(\mathbf{x}, \mathbf{y}) \geq \Delta_d(\mathbf{x}_d, \mathbf{y}_d),
    \end{equation*}
    where \( \mathbf{x}, \mathbf{y} \) are compositions with \( k \) parts, and \( \mathbf{x}_d, \mathbf{y}_d \) are corresponding subcompositions with \( d < k \) parts.
    \end{itemize}
    \end{alertblock}

    
      \begin{itemize}
        \item[1)] \textbf{Basic mathematical operations} on the simplex (perturbation, powering, inner product)
        \item[2)] \textbf{Measures of central tendency}: simplicial means and medians

        \item[3)] \textbf{Distances and metrics} between compositions
        \item[4)] \textbf{Distributions} defined on or adapted to the simplex
      \end{itemize}
}




\end{frame}

\begin{frame}{Frame Title}

     \note{

an example of principles, personaly after one year working on this topice I still dont know which problem this principle try to solve and I am not alone in this regard. there is a long-term contrvery around this topic.

}

\[
\mathbf{T}'_4 = T\Big(\frac{Y'_1}{S'_4},\frac{Y'_2}{S'_4},\frac{Y'_3}{S'_4},\frac{Y'_4}{S'_4}\Big):(T'_1,T'_2,T'_3,T'_4)
\\
\mathbf{T}'_3 = T\Big(\frac{Y'_1}{S'_4},\frac{Y'_2}{S'_4},\frac{Y'_3}{S'_4}\Big):(T'_1,T'_2,T'_3)
\\
D_3(\mathbf{T}_3,\mathbf{T}'_3) \le D_4(\mathbf{T}_4,\mathbf{T}'_4)
\]



\end{frame}

% \begin{frame}{A Fundamental Controversy}
%     There have been long term debates between researcher in this field from the beginning sofar.

% Aitchison (1982) pioneered efforts to address the analysis of compositional data by arguing that such data should be transformed to Euclidean space using log-ratio methods.
   
%   % \bibliographystyle{plainnat}
%   % \bibliography{references} 

    
% \end{frame}
% \begin{frame}{Key Critics: D.F.\ Watson}
% \note{Wantson and philip were the first one that chalenged this principles and from 1989 to 1992 depates exchanges between them and published in... journal}

%     \textbf{Watson \& Philip (1989)} were among the \textbf{earliest} voices
%     to challenge the log-ratio framework.
%  \only<1>{
%     \vspace{0.3cm}
%     \begin{block}{Core Argument}
%       \begin{itemize}
%         \item Questioned whether \textbf{scale invariance} is a
%               fundamental requirement for compositional analysis.
%         \item If compositions are strictly defined as summing to one,
%               scale invariance is \textbf{already built in} — and imposing
%               it as a separate principle becomes circular
%         \item Their exchanges with Aitchison (1989--1992) were
%               instrumental in \emph{sharpening} the CoDa Principles,
%               even if from an adversarial position.
%       \end{itemize}
%     \end{block}}
%     \only<2>{
% \note{
% Then Aitchison used the word "Stupid" in one of his note letter and undermind Watson knowledge... this mad Watson to leave the debates.

% Then Aitchison and other new researcher developed his idea and published more paper including a book...
% }
%     \begin{alertblock}{Historical Note}
%       The debate was cut short by a notable incident: Aitchison
%       used \textbf{derogatory language} toward Watson in the
%       published exchange, which effectively ended Watson's
%       participation.
 
    
%     \end{alertblock}
 
%     \vspace{0.3cm}
%     \begin{block}{Legacy}
%       Watson's questions about scale invariance remain
%       \textbf{unresolved} and were later taken up more formally
%       by Scealy \& Welsh (2011, 2014).
%     \end{block}

%  }
% \end{frame}

% \begin{frame}{Key Critics: Zier \& Rehder (1998, 2001)}
% \note{ Zier and Rehder in 1998 to 2001 started to chalenged CoDa School principles. One interesting point was even log-ration transformations are not satosfied principles in some sense.}
%     Zier \& Rehder continued the challenge against the CoDa School,
%     focusing on \textbf{technical and foundational} issues.
%  \only<1>{
%     \vspace{0.3cm}
%     \begin{block}{Core Arguments}
%       \begin{itemize}
%         \item Challenged the mathematical definition of a

%         \item \colortext{yellow}{Pointed out that even \textbf{Aitchison's own earlier
%               work} did not strictly adhere to the current CoDa
%               definition of scale invariance}
%               \textbf{metric} within compositional data analysis
%         \item Argued that scale invariance is \textbf{problematic}
%               when compositions are strictly constrained to sum to one
        
%         \item Made the case that log-ratio methods are not the
%               \emph{only} valid approach
%       \end{itemize}
%     \end{block}
% }
% \only<2>{
% \note{Despide they made some valid points but they most of thier points were ignored partly becuase of some technical errors that they made}
%     \begin{block}{Reception by the CoDa School}
%       Despite raising \textbf{some valid points}, their criticisms
%       were \textit{dismissed relatively easily} — partly because
%       of technical errors in their own work that the CoDa School
%       used to deflect the substantive arguments.
%     \end{block}
 
 
% }
 
% \end{frame}


\begin{frame}
\begin{itemize}
    \item[\textcolor{blue}{\large$\blacktriangleright$}] These transformations were introduced by Aitchison in 1982, but the principles of CoDA have been developing over the years since then.
    \item[\textcolor{blue}{\large$\blacktriangleright$}] In 1989, Watson and Philip challenged both the log-ratio approach and the notion that these principles should be treated as universal requirements, proposing alternatives such as directional approaches.
    \item[\textcolor{blue}{\large$\blacktriangleright$}] This sparked a debate between Watson and Aitchison that unfolded over several years through a series of published exchanges in \textit{Mathematical Geology}.
    \item[\textcolor{blue}{\large$\blacktriangleright$}] What began as a methodological disagreement gradually escalated in tone, becoming increasingly confrontational.
    \item[\textcolor{blue}{\large$\blacktriangleright$}] The debate ultimately ended with Watson's withdrawal --- Aitchison's use of the word \textcolor{red}{``stupid''} to dismiss Watson's contributions led Watson to disengage from further discussion.
\end{itemize}
\end{frame}

\begin{frame}
\begin{itemize}
    \item[1.] Following this period, the CoDA framework grew increasingly dominant:
    \begin{itemize}
        \item[1.1] Aitchison continued to advance the framework through numerous publications, including an influential book.
        \item[1.2] A broader research community developed around the field that supported these principles.
    \end{itemize}
    \item[2.] As a result, alternative approaches received considerably less attention.
\end{itemize}
\end{frame}


\begin{frame}{The Controversy: To Transform or Not?}
\only<1>{
The CoDA school has redefined these principles several times in response to criticism.



        \begin{alertblock}{ (Zier & Rehder (1998, 2001))}
     They pointed out that even Aitchison’s own earlier work did not strictly satisfy the later CoDA definition of scale invariance.
    \end{alertblock}


        \begin{alertblock}{ (Scealy \& Welsh 2014)}
      The CoDa Principles act as a \textbf{gatekeeping criterion}:
      valid and useful methods have been dismissed simply for not
      conforming to them.
 
      \vspace{0.2cm}
      This is \textbf{unusual, perhaps unique, in statistical
      data analysis}.
    \end{alertblock}
}
\end{frame}



\begin{frame}{One Idea: Transforming to Real Space}


\note{
Despite the strong emphasis of the CoDA school on log-ratio transformations, they are not the only way to work with compositional data.

Several other transformations have also been developed, including approaches that are not based on logarithms and do not strictly follow the classical CoDA framework.

Some of these alternatives are listed in the table.

And in practice, depending on the data and the modeling objective, they can sometimes perform as well as — or even better than — log-ratio transformations.


Still, even when these transformations are useful, one important difficulty remains.

Once we move the data out of the simplex, interpretation becomes less direct.

Some results, such as predictions, can often be mapped back into the simplex.

But many others become much harder to interpret in the original compositional space.

So while transformation can be useful computationally, it often makes interpretation less natural.

This is one of the main reasons why working directly in the simplex can be valuable.

And this is exactly the direction we followed in the cotext of machine learning:


}

\only<1>{
  \framesubtitle{A transformation $\mathcal{S}^{K-1} \to \mathbb{R}^{K-1}$ unlocks standard inference: regression, PCA, LDA}

  \vspace{0.2cm}
  \renewcommand{\arraystretch}{1.35}
  \setlength{\tabcolsep}{8pt}

  \begin{tabular}{>{\bfseries}l l >{\footnotesize\itshape\color{gray}}l}
    \toprule
    \rowcolor{mainblue}
    \color{white}Method &
    \color{white}Formula &
    \color{white}Reference \\
    \midrule

    \rowcolor{accentblue!20}
    ALR & $\log y_i - \log y_K$ & Aitchison (1982) \\

    \rowcolor{accentblue!20}
    CLR & $\log y_i - K^{-1}\sum_{j=1}^{K} \log y_j$ & Aitchison (1983) \\

    \rowcolor{accentblue!20}
    ILR & $\mathbf{H}\!\left(\log y_i - K^{-1}\sum_{j=1}^{K} \log y_j\right)$ & Egozcue et al.\ (2003) \\

    \midrule
  
    {\normalfont $\alpha$} &
    {\footnotesize $\mathbf{H}[K_\alpha^{-1}(\mathbf{y} \otimes \alpha)]$} &
    Tsagris et al.\ (2011) \\

    {\normalfont BC$_1$} &
    {\footnotesize $\lambda^{-1}\!\left[(y_i / y_K)^\lambda - 1\right]$} &
    Aitchison (2003) \\

    {\normalfont BC$_2$} &
    {\footnotesize $\theta^{-1}(y_i^\theta - 1)$} &
    Greenacre (2010) \\

    {\normalfont Hyper-spherical} &
    {\footnotesize $\arccos\!\left(\dfrac{y_i}{\prod_{j=i+1}^{K} y_j}\right)$} &
    Wang et al.\ (2007) \\

    {\normalfont Square root} &
    {\footnotesize $\sqrt{y_i}$} &
    Stephens (1982) \\

    \bottomrule
  \end{tabular}

  \vspace{0.2cm}
  {\small\textbf{Key:} ALR, CLR, ILR are the most widely used;
  all resolve singularity of $\Sigma_{\mathbf{X}}$ by mapping to unconstrained $\mathbb{R}^{K-1}$.}
}



\end{frame}


